\documentclass[11pt,reqno,a4paper]{amsart}
\usepackage[T1]{fontenc}
\usepackage{amssymb,booktabs}
\usepackage[hidelinks]{hyperref}
\allowdisplaybreaks[2]
\raggedbottom
\setlength{\emergencystretch}{1.8em}

\newtheorem{theorem}{Theorem}[section]
\newtheorem{proposition}[theorem]{Proposition}
\newtheorem{corollary}[theorem]{Corollary}
\newtheorem{lemma}[theorem]{Lemma}
\theoremstyle{remark}
\newtheorem{remark}[theorem]{Remark}
\newtheorem*{theoremT}{Theorem T}

\newcommand{\Fac}{\operatorname{Fac}}
\newcommand{\e}{\varepsilon}
\newcommand{\Y}{\mathcal Y}
\newcommand{\tern}[1]{\mathtt{#1}}
\newcommand{\R}{\mathbb R}
\newcommand{\Q}{\mathbb Q}
\newcommand{\Z}{\mathbb Z}
\newcommand{\dist}[1]{\lVert#1\rVert}
\newcommand{\rev}{\mathrm{rev}}
\newcommand{\word}[1]{\mathtt{#1}}
\newcommand{\HH}{\mathcal H}

\title[A rational translation of the middle-third Cantor set]
{A rational translation bounding palindromes on the middle-third Cantor set}
\newsavebox{\affilbox}
\address{School of Artificial Intelligence, Shanghai Jiao Tong University, Shanghai, China}
\email{mqyqwerty@sjtu.edu.cn}
\date{}
\subjclass[2020]{11J68, 11J83, 68R15}
\keywords{Middle-third Cantor set, algebraic irrationals, ternary expansions, palindromes, rational translations, Mahler's problem}
\hypersetup{pdfauthor={Qiyuan Ma},pdftitle={A rational translation bounding palindromes on the middle-third Cantor set},pdfdisplaydoctitle=true}

\begin{document}
\sbox{\affilbox}{%
  \begin{minipage}{0.92\textwidth}
    \centering\normalfont\normalsize
    School of Artificial Intelligence,\\
    Shanghai Jiao Tong University, Shanghai, China\\
    ORCID: \href{https://orcid.org/0009-0000-9532-0821}{0009-0000-9532-0821}\\
    \href{mailto:mqyqwerty@sjtu.edu.cn}{mqyqwerty@sjtu.edu.cn}
  \end{minipage}}
\author[Qiyuan Ma]{Qiyuan Ma\\[6pt]\protect\usebox{\affilbox}}

\begin{abstract}
Let $C_3$ be the middle-third Cantor set. We construct a single rational number $\tau$ such that, for every irrational $x\in C_3$, the ternary expansion of $x+\tau$ has palindrome length at most $145$. A direct construction, which does not rely on a finite search, gives the weaker bound $847$. Consequently, if every real algebraic irrational has ternary palindromes of unbounded length, then $C_3$ contains no algebraic irrational. The latter assertion is Mahler's problem, which remains open. We also record a circle-return criterion for the occurrence of a square, and a brief entropy remark: a base-$b$ expansion whose factor entropy is strictly less than $\tfrac12\log b$ admits a rational translation with bounded palindromes. The two unbounded-length questions in Bugeaud's Problem~10.42 remain open, as does the existence of an algebraic irrational in $C_3$.
\end{abstract}
\maketitle

\section{Introduction}\label{sec:intro}

For an integer $b\ge2$ and a real number $x$, write $\{x\}=x-\lfloor x\rfloor$ and $\dist{x}=\min_{n\in\Z}|x-n|$. Let $\mathsf S_b(x)$ mean that the canonical base-$b$ expansion of the fractional part of $x$ contains factors $VV$ with $|V|$ unbounded. Let $\mathsf P_b(x)$ mean that it contains palindromes of unbounded length. Bugeaud's Problem~10.42 asks whether each property holds for every real algebraic irrational $x$ and every integer $b\ge2$ \cite[p.~220]{Bugeaud}. Adamczewski and Rampersad posed the corresponding questions in base two \cite{AR}.

The digit block starting after $i$ places and of length $d$ has integer value
\begin{equation}\label{eq:digitblock}
 D_{i,d}(x)=\lfloor b^{i+d}x\rfloor-b^d\lfloor b^i x\rfloor.
\end{equation}
Leading zeros are included. For a word $W$ let $[W]_b$ denote its integer value. Deleting any fixed prefix preserves both $\mathsf S_b$ and $\mathsf P_b$. A long palindrome can be trimmed symmetrically. If all unbounded squares began in a finite set of starts, one start would have unbounded roots. The second half of a sufficiently large square at that start contains each earlier given square at the same start, and moves that square beyond the deleted segment.

The long-repetition criterion of Adamczewski--Bugeaud--Luca, as stated by Bilu \cite[Theorem~3.3, p.~8]{Bilu}, implies that the greatest length of two disjoint equal factors among the first $N$ digits of a fixed real algebraic irrational is $o(N)$. The corresponding length $L_b(N)$ with distinct starts and possible overlap also satisfies
\begin{equation}\label{eq:known-early}
 L_b(N)=o(N).
\end{equation}
Indeed, overlapping matches of length $\ell$ at distance $d<\ell$ give a period-$d$ segment of length $\ell+d$. It contains $\lfloor\ell/d\rfloor+1$ adjacent copies of a length-$d$ word. Splitting them into two equal groups gives a disjoint match of length at least $\ell/3$, as in \cite[Lemma~3.4]{Bilu}.

Let $C_3$ be the classical middle-third Cantor set. Every point of $C_3$ admits a ternary representation with digits $0$ and $2$. Mahler's problem asks whether $C_3$ contains any algebraic irrational; the question remains open, including for quadratic irrationals.

The main result of this paper is a uniform rational translation that bounds palindrome lengths on the whole irrational part of $C_3$. Theorem~\ref{thm:cantor} produces one rational $\tau$ such that every irrational $x\in C_3$ has ternary palindrome length at most $145$ in $x+\tau$. The same argument yields a direct, non-search construction with the weaker bound $847$. Rational translation preserves algebraicity and irrationality. Therefore a positive answer to $\mathsf P_3$ for all algebraic irrationals would imply that $C_3$ contains no algebraic irrational. This is a relation between two open problems, not a solution of either of them.

The first subsection records a circle-return criterion: a sufficiently close return of $(b^{n+h}-b^n)x$ forces a square of root length $h$ in some rational translate, and, for expansions supported on a sufficiently separated digit set, the existence of unbounded square roots is equivalent to a quantitative return condition. For the middle-third Cantor set in base three the threshold is $3^{-h}$.

A closing remark observes that if a base-$b$ word has factor entropy strictly less than $\tfrac12\log b$, then a random periodic addend produces, with positive probability, a rational translation with bounded palindromes. In particular, if the palindrome assertion of Problem~10.42 were true in base $b$ for all algebraic irrationals, every such number would have factor entropy at least $\tfrac12\log b$. This entropy comparison is developed more fully elsewhere; we include only the short argument that sits naturally with the translation construction.

Both unbounded-length questions of Problem~10.42 remain open, and so does Mahler's problem.

\section{Exact digit comparisons and rational translations}\label{sec:geometric}

\subsection{When a circle return is a square}
Let $B=b^h$, $X=\{b^nx\}$, and $Y=\{BX\}$. A square of root length $h$ starting after $n$ digits occurs precisely when $X$ and $Y$ belong to a common half-open cell $[a/B,(a+1)/B)$. For example, if $X=(a+v)/B$ with $0\le v<1$, then $Y=v$, so the second block has value $a$ exactly when $a/B\le v<(a+1)/B$. This cell condition gives both of the following consequences.

\begin{proposition}\label{prop:circle}
Let $q\ge2$ be coprime to $b$ and $g=\gcd(q,b^h-1)\ge2$. If
\[
 \dist{(b^{n+h}-b^n)x}< (1-1/g)b^{-h},
\]
then some $x+k/q$, with $0\le k<q$, has a root-$h$ square starting after $n$ digits.

Suppose instead that an irrational $x$ has a tail over a digit set $D$ of size at least two. Put
\[
 \omega=\frac{\max D-\min D}{b-1},
 \qquad j_D=\min_{\substack{a,a'\in D\\a\ne a'}}|a-a'|.
\]
If $j_D\ge2\omega$, then $\mathsf S_b(x)$ is equivalent to the existence, for unbounded $h$, of some $n$ in that tail with
\[
 \dist{(b^{n+h}-b^n)x}<(j_D-\omega)b^{-h}.
\]
For $D=\{0,2\}$ and $b=3$ the threshold is $3^{-h}$.
\end{proposition}

\begin{proof}
For the first assertion, take only $k$ divisible by $q/g$. Each such translation rotates both $X,Y$ by the same angle, and the possible angles form the equally spaced $g$-grid. Multiplying the circle cut phases by $B$ still permutes the grid because $\gcd(B,g)=1$. The shorter arc between $X$ and $Y$ has length less than $(1-1/g)/B$, so after scaling by $B$ it cannot cover every cut phase. For a missed phase, $X,Y$ land in one cell. If the two points coincide, all rotations work.

For the second assertion, distinct length-$h$ integers over $D$ differ by at least $j_D$. This follows by taking the first differing digit and bounding all subsequent digits by the full base range. The remainder tails range over an interval of width $\omega$. In ordinary distance, unequal blocks are therefore at least $(j_D-\omega)/B$ apart. If $\omega<1$, a short arc through zero is ruled out for all sufficiently large $h$. If $\omega=1$, a return through zero below the threshold forces the first blocks to be the two extreme constant blocks. Indeed a nonextreme block is at least $j_D/B$ from the relevant endpoint even before adding the nonnegative distance from the other point. The constant words of length $h$ then contain squares of root length $\lfloor h/2\rfloor$. Conversely, matching first blocks leave a distance strictly less than $\omega/B$, since an irrational tail never attains either endpoint. This gives the inequality as $j_D\ge2\omega$.
\end{proof}

The finite-grid constant is sharp along any unbounded sequence of $h$ with fixed $g\ge2$. In fact, take $d=(1-1/g+\epsilon)/B$, with $0<\epsilon<1/g$, and $X_i=(i+d)/(B-1)$ for $0\le i<B-1$. Then $\{BX_i\}=\{X_i+d\}$. The phases of $BX_i$ form a $(B-1)$-grid up to a fixed rotation. For large $B$, choose a phase so the scaled arc contains every cut phase. The nonaligned translations have relative rotation at least $1/q$, so cannot give one cell for large $B$. These examples are at individual finite scales.

\subsection{A rational translation of the middle-third Cantor set}
Let $C_3$ be the classical middle-third Cantor set. Every point admits a ternary representation with digits $0,2$. For the arithmetic consequence it suffices to work with irrational points, which have no endpoint ambiguity.

\begin{theorem}\label{thm:cantor}
There is one rational $\tau$ such that all irrational $x\in C_3$ have ternary palindrome length at most $145$ in $x+\tau$. There is also a direct, non-search construction with bound $847$. In particular, the positive answer to $\mathsf P_3$ for all algebraic irrationals implies the conjecture that $C_3$ contains no algebraic irrational.
\end{theorem}

\begin{proof}
Use (ordinary) addition in base three. At each digit let $u$ be the digit of the addend, $a\in\{0,2\}$ the input and $c\in\{0,1\}$ the incoming carry. The output is $a+u+c$ modulo $3$, with outgoing carry $\lfloor(a+u+c)/3\rfloor$. Scanning a palindrome inward uses a pair state $(L,R)$, consisting of the outgoing carry at the left column and the incoming carry at the right column. A transition to $(L',R')$ requires equal output digits from the columns with input carries $L',R$. This checks all possible actual left and right carries.

Use column pairs $A=(0,2)$, $B=(2,0)$, $C=(0,1)$. Directly applying the digit equation, the sequence $A,B,C,A,B$ takes the full state set through
\[
 \{00,01,10,11\}\longrightarrow\{00,01,11\}
 \longrightarrow\{00,11\}\longrightarrow\{01,10\}
 \longrightarrow\{01\}\longrightarrow\varnothing .
\]
Thus the left word $U=02002$ and the reflected right word $V^{\rm rev}=02102$ cannot match, whatever the carries. Pack this pair in a periodic addend as follows. Write $\ell=5$, $P=\ell(\ell+1)$, $W=2(P+\ell)$, and use six macroblocks of width $W$ with colors $001011$. In color zero take the first $W$ letters of $U^\omega$. In color one take the first $W$ letters of $(V^{\rm rev}0)^\omega$. Every reflection of the six colors pairs unequal colors. For a reflection of the whole period, write the reflection sum as $aW+r$. On the two coordinate ranges $i\le r$ and $i>r$, the macroblock reflection sums are $a$ and $a-1$ respectively. The longer range has length at least $P+\ell$. It therefore contains an entire occurrence of $V^{\rm rev}$ starting at a prescribed residue modulo $\ell$, since the possible starts advance by $\ell+1$ and run through every residue over $P$ positions. Choose the residue so its reflected start is a multiple of $\ell$ in the opposite macroblock. This gives the bad pair at a phase for every reflection. A left half of length at least $6W+\ell-1$ contains the necessary full block at that phase. Consequently one cannot have a palindrome of length at least $2(6W+\ell-1)$. This is the bound $847$.

For the smaller constant, let $T$ be the concatenation of the following two rows and use $\tau=[T]_3/(3^{142}-1)$.
\begin{center}
\footnotesize\ttfamily
211020102120122112001122021201202221201220122210201021112002210211021110\\
2011202022220101212211222020202211212021211020021002202120202022211000
\end{center}
For each of the $142$ phases and each parity, start from all four states and iterate the column rule. For the even test use $62$ columns on each side and require equal carries at the center. For the odd test use $73$ on each side and require an allowed central digit. The verifier reproduces all $284$ rejections. Hence there is no even palindrome of length $124$ and no odd one of length $147$. Taking central factors gives the stated bound. Rational translation preserves algebraicity and irrationality, so any algebraic irrational in $C_3$ produces a counterexample to the ternary palindrome assertion.
\end{proof}

The same five-column calculation applies to any nonadjacent digit set containing the two extreme digits, with $2$ replaced by $b-1$. For a nonadjacent set omitting the top digit after normalization, the column sequence $(0,0),(0,0),(b-1,0)$ suffices. An affine normalization followed by doubling also covers a digit interval of half the full width. The supplemental carry graphs handle some larger digit sets conditionally upon square avoidance. Their role is to force a periodic input along any sufficiently long attempted palindrome path.

One further consequence concerns entropy. If a base-$b$ word has factor count $p(n)$ and $\lim_n n^{-1}\log p(n)<\frac12\log b$, choose a random length-$n$ periodic addend. For each fixed factor of length $n$ or $n-1$ and each phase, the shifted block is uniform for either prescribed incoming carry. The number of palindromes of length $j$ is $b^{\lceil j/2\rceil}$. The union cost for both lengths is at most
\[
 2n\bigl(p(n)b^{-\lfloor n/2\rfloor}+p(n-1)b^{-\lfloor(n-1)/2\rfloor}\bigr),
\]
which tends to zero. Some rational translation thus has bounded palindromes. If the palindrome assertion were true in base $b$ for all algebraic irrationals, each such number would have factor entropy at least $\frac12\log b$.


\section*{Numbering correspondence}
For ease of cross-reference, the following table records the numbering of the principal statements as they appeared in an unpublished extended account of related results dated 26~November~2026.
\begin{center}
\begin{tabular}{@{}ll@{}}
\toprule
This paper & Extended account \\
\midrule
Proposition~\ref{prop:circle} & Proposition~3.1 \\
Theorem~\ref{thm:cantor} & Theorem~3.2 \\
\bottomrule
\end{tabular}
\end{center}

\begin{thebibliography}{99}
\bibitem{AR}
B. Adamczewski and N. Rampersad,
\emph{On patterns occurring in binary algebraic numbers},
Proc. Amer. Math. Soc. 136 (2008), 3105--3109.
\href{https://doi.org/10.1090/S0002-9939-08-09319-2}{DOI}.
\href{https://adamczewski.perso.math.cnrs.fr/Overlap.pdf}{Author version}.

\bibitem{Bugeaud}
Y. Bugeaud,
\emph{Distribution Modulo One and Diophantine Approximation},
Cambridge Tracts in Mathematics, vol.~193, Cambridge University Press,
2012.
\href{https://doi.org/10.1017/CBO9781139017732}{DOI}.

\bibitem{Bilu}
Y. F. Bilu, \emph{The many faces of the subspace theorem [after
Adamczewski, Bugeaud, Corvaja, Zannier...]}, S\'eminaire Bourbaki,
expos\'e 967, Ast\'erisque 317 (2008), 1--38.
\href{https://www.numdam.org/book-part/AST_2008__317__1_0/}{Publication record}.
\href{https://arxiv.org/pdf/0907.2098#page=8}{Theorem 3.3, printed p.~8}.


\end{thebibliography}

\end{document}
